% TT09032C
% Rückenmarksverletzung mit Symptomen
% 14.0
% 2013-11-30
:ref:`m-ws-trauma`
Taktik \Ca{Zeitkritisch! Vitale Bedrohung!} Zügige Bergung,
Immobilisation und Transport in eine geeignete Einrichtung}


Vgl. \prettyref{m:einschaetzung-indikation-ws-immobilisation}
\begin{itemize}
  -   Bei ansprechbaren Patienten: immer
  \EE{\Abk{HWS}-Immobilisation} (HWS-Schiene)
  -   **Bergung** mit Schaufeltrage, Spineboard oder
  KED-System
  -   **Lagerung** und Schienung mittels
  Vakuummatratze, Schaufeltrage mit Fixationsmöglichkeit oder
  Spineboard.\footnote{Welches Hilfsmittel
  eingesetzt wird ist abhängig vom Patientenzustand,
  dem Training des Teams, internen Vorschriften und
  vorhandenem Material.}
  -   Bei Querschnittslähmung: Kennzeichnung der
  Sensibilitätsgrenze mit Angabe der Uhrzeit
  -   **Monitoring**: Achte besonders auf Störungen
  von Atmung und Bewußtsein, sowie Schockentwicklung.
%   -   **Unfallmechanismus** erheben und dokumentieren 
  -   **Transportentscheidung**: Schockraum
\end{itemize}